
\section{Preliminaries}

\subsection{Statistical watermarking}

Statistical watermarking for Large Language Models (LLMs) provides a probabilistic framework for distinguishing machine-generated text from human-written text. Unlike post-hoc detection methods that rely on classifiers trained on model artifacts, watermarking actively embeds a statistical signal into the generation process. This signal is designed to be imperceptible to humans yet statistically significant to an algorithmic detector possessing a secret key.

\paragraph{Statistical watermarking as a hypothesis testing problem.}
Formally, statistical watermarking is cast as a hypothesis testing problem. Let $q$ denote the target distribution over the sample space $\mathcal{V}$, and let $\mathcal{S}$ represent the space of pseudorandom seeds. The watermarking protocol modifies the standard decoding mechanism of the LLM such that the output tokens $V \in \mathcal{V}$ and the pseudorandom seeds $S \in \mathcal{S}$ are sampled jointly from a watermarked distribution $\mathcal{P}_W$. 

The detection task seeks to determine if a given observed text $v$ was generated by the watermarked model: $\mathcal{D}: \mathcal{V} \times \mathcal{S} \to \{\text{Watermarked}, \text{Unwatermarked}\}$. This reduces to testing for independence between the tokens $v$ and the seeds $s$ reconstructed via the detector's key. We define the null and alternative hypotheses as follows~\citep{huang2023towards}:
\begin{itemize}
    \item \textbf{Null hypothesis $\textbf{H}_0$:} The text $v$ is generated independently of the seeds $s$ (e.g., by a human or an unwatermarked model).
    \item \textbf{Alternative hypothesis $\textbf{H}_1$:} The text $v$ and seeds $s$ are sampled from the joint watermarked distribution $\mathcal{P}_W$, implying they come from the watermarked model.
\end{itemize}

It is pertinent to distinguish this \textit{distributional} watermarking approach from classical instance-level watermarking techniques applied to static media, such as images or audio \cite{cox2007digital}. While classical methods embed a signature into a specific, fixed outcome (post-generation), statistical watermarking modifies the stochastic sampling process itself, ensuring that any realization from the model carries the statistical evidence of its origin without requiring rigid alterations to the final output.

\paragraph{Statistical guarantees.}
A robust watermarking scheme is characterized by its ability to provide three fundamental guarantees.

First, to preserve generation quality and ensure watermarked content remains indistinguishable from ordinary samples, the distance (e.g., KL-divergence) between the watermarked distribution $\mathcal{P}_W$ and the original target distribution $q$ must be minimized. Ideally, we seek a scheme that introduces zero distributional shift~\citep{hu2023unbiased}.

\begin{definition}[Distortion-free]\label{def:unbias}
    A watermark is considered \textit{distortion-free} (or unbiased) if the outcome marginal of the watermarked distribution $\mathcal{P}_W$ matches the target distribution $q$. Formally, a distortion-free watermark satisfies:
    \begin{align*}
        \sum_{s \in \mathcal{S}} \mathcal{P}_W(A, s) = q(A), \quad \forall A \subseteq \mathcal{V}.
    \end{align*}
\end{definition}

Second, it is often required in practice that the detector operates without access to the watermarked distribution since the underlying model's parameters are generally unknown to the detector. This means that the detection scheme needs to be designed for a \emph{family of target distributions} in a model-agnostic way~\citep{kuditipudi2023robust}, leading to the next concept.

\begin{definition}[Model-agnosticity]\label{def:agnostic}
    A watermark is \textit{model-agnostic} if the seed-marginal of $\mathcal{P}_W$ is chosen independently of the target distribution $q$. That is, the distribution of seeds $S$ does not depend on the model's logits.
\end{definition}

Third, the reliability of the watermark is quantified by standard statistical error rates.
The \emph{Type I error} (false positive rate) measures the probability of incorrectly identifying non-watermarked text (e.g., human-written) as watermarked. Given the high stakes of false accusations (e.g., in academic integrity), this error must be bounded by a small significance level $\alpha$:
\begin{align*}
    \mathbb{P}_{H_0}(\mathcal{D}(V, S) = \text{Watermarked}) \leq \alpha.
\end{align*}
The \emph{Type II error} (false negative rate) measures the probability that watermarked text fails to be detected. Minimizing this error is equivalent to maximizing the statistical power of the test. For a target error rate $\beta$, we require:
\begin{align*}
    \mathbb{P}_{H_1}(\mathcal{D}(V, S) = \text{Unwatermarked}) \leq \beta.
\end{align*}


\subsection{Sequential hypothesis testing and e-values}

Consider a measurable space $(\Omega, \mathcal{F})$ and a family of probability distributions $\mathcal{P}$. We are interested in testing a null hypothesis $\mathcal{H}_0: P \in \mathcal{P}_0 \subset \mathcal{P}$ against an alternative $\mathcal{H}_1: P \in \mathcal{P}_1 = \mathcal{P} \setminus \mathcal{P}_0$. 
An e-value is a random variable whose expectation is bounded by unity under the null hypothesis~\citep{vovk2021values}.

\begin{definition}[E-value]
A nonnegative random variable $E$ is an e-value for $\mathcal{H}_0$ if for all $P \in \mathcal{P}_0$,
\begin{align*}
    \mathbb{E}_P[E] \le 1.
\end{align*}
\end{definition}

Intuitively, an e-value represents the amount of wealth a gambler would have after betting against the null hypothesis, in a game where the game is fair or unfavorable under $\mathcal{H}_0$, and starting with an initial wealth of one. On the other hand, if $E$ becomes large, this suggests that the null hypothesis is unlikely to be true.

\paragraph{Sequential evidence and stopping time guarantees.}

The primary advantage of e-values lies in their behavior regarding stopping times. In a sequential setting, we are given a filtration $(\mathcal{F}_t)_{t \ge 0}$ and construct a sequence of e-values $(E_t)_{t \ge 0}$. If $(E_t)_{t \ge 0}$ is a nonnegative supermartingale with respect to the null distributions (often called a \textit{test martingale}), it allows for \textit{anytime-valid} inference.
This property derives from Ville's inequality, a time-uniform generalization of Markov's inequality.

\begin{theorem}[Ville's inequality]
Let $(E_t)_{t \ge 0}$ be a nonnegative supermartingale with respect to $P \in \mathcal{P}_0$ such that $E_0 \le 1$. Then for any $\alpha \in (0,1)$,
\begin{align*}
    P\left( \exists t \ge 0 : E_t \ge \frac{1}{\alpha} \right) \le \alpha.
\end{align*}
\end{theorem}

This result guarantees that a researcher can track the e-value process continuously and stop at any data-dependent time $\tau$ (a stopping time) while maintaining Type-I error control. Specifically, $\mathbb{E}_P[E_\tau] \le 1$ holds for any stopping time $\tau$ (possibly unbounded), a kind of guarantee which is not available for p-values.
